\section{Resultados}\label{sec:resultados}

\subsection{Métricas de implementação e catálogo de ferramentas}

A implementação reúne 42 arquivos TypeScript (14.677 linhas) na camada MCP e
54.442 linhas em Python na camada de interface, das quais 38.050 no módulo de
comandos e 6.529 no roteador principal (\texttt{kicad\_interface.py}). O
servidor registra \textbf{233 ferramentas}, organizadas em 24 módulos
funcionais, agrupadas na Tabela~\ref{tab:categorias} em dez macro-áreas.

\begin{table}[htbp]
\centering
\small
\caption{Ferramentas MCP por macro-área funcional.}
\label{tab:categorias}
\begin{tabular}{@{}p{5.6cm}p{1.6cm}p{7.2cm}@{}}
\toprule
\textbf{Macro-área} & \textbf{N.\ de ferramentas} & \textbf{Módulos de origem} \\
\midrule
Esquemático (autoria, lote, hierarquia, acabamento) & 65 & \texttt{schematic.ts}, \texttt{schematic-batch.ts}, \texttt{schematic-hierarchy.ts}, \texttt{schematic-layout.ts} \\
Bibliotecas e criação de partes & 31 & \texttt{library.ts}, \texttt{library-symbol.ts}, \texttt{footprint.ts}, \texttt{symbol-creator.ts} \\
Projeto e ciclo de vida da placa & 31 & \texttt{project.ts}, \texttt{board.ts} \\
Componentes & 28 & \texttt{component.ts} \\
Exportação e fabricação & 27 & \texttt{export.ts} \\
Roteamento e autorouter & 21 & \texttt{routing.ts}, \texttt{freerouting.ts} \\
Datasheets e cadeia de suprimentos & 13 & \texttt{datasheet.ts}, \texttt{jlcpcb-api.ts}, \texttt{parts-registry.ts}, \texttt{digikey-api.ts} \\
Regras de projeto e validação & 9 & \texttt{design-rules.ts}, \texttt{validation.ts} \\
Interface do KiCad e descoberta & 6 & \texttt{ui.ts}, \texttt{router.ts} \\
Importação de formatos legados & 2 & \texttt{eagle.ts}, \texttt{pcb-import.ts} \\
\midrule
\textbf{Total} & \textbf{233} & 24 módulos \\
\bottomrule
\end{tabular}
\end{table}

Do total, 173 ferramentas (74,2\,\%) estão indexadas para descoberta por
palavra-chave em 16 categorias, conforme a Equação~\ref{eq:cobertura}; as 60
restantes permanecem executáveis, porém não localizáveis pelos mecanismos de
busca --- número que um teste automatizado congela para que só possa diminuir.
Além das ferramentas, o servidor expõe 23 recursos dinâmicos com o estado do
projeto (placa aberta, componentes, bibliotecas e informações do projeto) e
modelos de \emph{prompt} para tarefas recorrentes de projeto.

\begin{equation}
\label{eq:cobertura}
C_{desc} = \frac{173}{233} \approx 74{,}2\,\%.
\end{equation}

\subsection{Execução das suítes de teste}

A Tabela~\ref{tab:testes} consolida a execução local de 10 de setembro de 2026.
No total, 2.451 casos foram coletados; 2.412 foram aprovados (98,4\,\%), com 8
falhas e 31 casos ignorados. A suíte TypeScript passou integralmente
(84 de 84).

\begin{table}[htbp]
\centering
\small
\caption{Resultados da execução local das suítes de teste (10/09/2026).}
\label{tab:testes}
\begin{tabular}{@{}p{4.0cm}rrrrp{2.4cm}@{}}
\toprule
\textbf{Suíte} & \textbf{Casos} & \textbf{Aprov.} & \textbf{Falhas} & \textbf{Ignor.} & \textbf{Tempo} \\
\midrule
TypeScript (Vitest 2.1) & 84 & 84 & 0 & 0 & 1,5\,s \\
Python (pytest 9.1) & 2.367 & 2.328 & 8 & 31 & 45,0\,s \\
\midrule
\textbf{Total} & \textbf{2.451} & \textbf{2.412} & \textbf{8} & \textbf{31} & 46,5\,s \\
\bottomrule
\end{tabular}
\end{table}

As oito falhas concentram-se em casos que invocam o utilitário
\texttt{kicad-cli} como processo externo (importação Eagle, exportação de
netlist, verificação ERC e compatibilidade de formato de esquemático) e em um
teste de estado de módulo legado. A análise dos rastros de execução indica
sensibilidade ao ambiente de validação --- instalação do KiCad por pacote
\emph{snap}, sem o diretório \texttt{/usr/share/kicad/schemas} esperado por
alguns utilitários --- e não a regressões funcionais nas ferramentas exercitadas
pelos demais 2.328 casos aprovados. Os registros completos das execuções estão
preservados no repositório (\texttt{evidencias/}).

Em integração contínua, o projeto executa a suíte TypeScript em quatro
sistemas operacionais (Ubuntu 24.04, Ubuntu 22.04, Windows e macOS) com Node.js
20 e 22, a suíte Python em Ubuntu com Python 3.9 a 3.12, e dois portões de
qualidade adicionais: o \emph{lint} TypeScript e a verificação de que o
inventário de ferramentas está em dia com o código.

\subsection{Estudo de caso: placa seguidora de linhas de quatro camadas}

O projeto analisado possui 97 \emph{footprints} distribuídos em quatro camadas
de cobre (F.Cu, In1.Cu, In2.Cu e B.Cu), espessura total de 1,6\,mm e contorno
de aproximadamente 82,1\,mm $\times$ 80,2\,mm. O esquemático é hierárquico,
com nove folhas (raiz, fonte regulada, \emph{MPPT}, pontes H, sensores ópticos
reflexivos, interface USB-C e demais blocos), e a lista de materiais agrupa 37
itens. As Figura~\ref{fig:render-topo} e~\ref{fig:render-fundo} mostram
renderizações 3D geradas a partir do arquivo de placa; a
Figura~\ref{fig:esquematico} apresenta a folha raiz do esquemático, cujo bloco
de título registra a autoria do projeto.

\begin{figure}[htbp]
\centering
\includegraphics[width=0.82\textwidth]{pcb-render-topo.png}
\caption{Vista superior da placa seguidora de linhas (renderização 3D gerada
com \texttt{kicad-cli pcb render}).}
\label{fig:render-topo}
\end{figure}

\begin{figure}[htbp]
\centering
\includegraphics[width=0.82\textwidth]{pcb-render-fundo.png}
\caption{Vista inferior da placa, evidenciando o plano de cobre e a
distribuição dos componentes SMD.}
\label{fig:render-fundo}
\end{figure}

\begin{figure}[htbp]
\centering
\includegraphics[width=0.95\textwidth]{esquematico-raiz.png}
\caption{Folha raiz do esquemático hierárquico (nove folhas), exportada com
\texttt{kicad-cli sch export pdf}.}
\label{fig:esquematico}
\end{figure}

A verificação de regras de projeto com \texttt{kicad-cli} 9.0.7 reportou
11 violações: 9 erros e 2 avisos, detalhados na
Tabela~\ref{tab:drc}. Não foram encontrados itens não conectados
(\emph{unconnected items}) nem divergências de paridade entre placa e
esquemático.

\begin{table}[htbp]
\centering
\small
\caption{Violações de DRC reportadas por \texttt{kicad-cli} 9.0.7 em 10/09/2026.}
\label{tab:drc}
\begin{tabular}{@{}p{3.4cm}p{1.9cm}p{1.2cm}p{8.0cm}@{}}
\toprule
\textbf{Tipo} & \textbf{Severidade} & \textbf{Qtde.} & \textbf{Diagnóstico} \\
\midrule
\texttt{annular\_width} & erro & 2 & Largura anular abaixo do mínimo de 0,100\,mm (medição de $-0{,}003$\,mm). \\
\texttt{clearance} & erro & 2 & Afastamento de 0,1812\,mm contra 0,200\,mm da classe \emph{Default}. \\
\texttt{malformed\_courtyard} & erro & 1 & \emph{Footprint} com \emph{courtyard} autointersectante. \\
\texttt{starved\_thermal} & erro & 4 & Conexão térmica à zona incompleta (1 de 2 raios exigidos). \\
\texttt{padstack} & aviso & 2 & \emph{Padstack} questionável (furo PTH sem cobre remanescente). \\
\bottomrule
\end{tabular}
\end{table}

O fluxo de projeto também produziu artefatos de fabricação no diretório
\texttt{production/} (lista de materiais, designadores, netlist IPC, arquivo de
posições e pacote compactado da placa), demonstrando a integração entre a
edição do projeto e a preparação para manufatura.

\subsection{Integrações com o ecossistema de fabricação}

O servidor integra três fontes de dados de componentes: a base local e a API
da JLCPCB~\citep{jlcpcb} (mais de 2,5 milhões de registros catalogados), o
enriquecimento de \emph{datasheets} via LCSC~\citep{lcsc} e a API
\emph{Product Information V4} da DigiKey~\citep{digikey}, utilizada tanto para
busca pontual de componentes quanto para varredura de bibliotecas de símbolos
em busca de itens obsoletos ou fora de estoque. O roteador automático
Freerouting~\citep{freerouting} é suportado via Java, Docker ou Podman. As
integrações com APIs externas dependem de credenciais opcionais definidas no
ambiente e degradam de forma controlada quando ausentes.
